\chapter{Lezione 15 --- Dead-beat control, sistemi stabilizzabili, stato aumentato e introduzione all'osservatore}

\section{Dead-beat controller}

La prima osservazione della lezione riguarda un sistema a tempo discreto
\[
x(k+1)=Ax(k)+Bu(k),
\]
controllato con una retroazione completa dello stato
\[
u(k)=-Kx(k).
\]

La dinamica in ciclo chiuso è quindi
\[
x(k+1)=(A-BK)x(k).
\]

Nel tempo discreto è possibile scegliere gli autovalori della matrice \(A-BK\) \emph{tutti in zero}.
Questo porta al cosiddetto \textbf{dead-beat controller}.

\subsection{Idea fondamentale}

Se riesco ad assegnare
\[
\lambda_1=\lambda_2=\cdots=\lambda_n=0,
\]
allora il polinomio caratteristico del sistema chiuso è
\[
p_{cl}(z)=z^n.
\]

In questo caso la matrice \(A-BK\) è nilpotente e, per un sistema di ordine \(n\),
lo stato va esattamente a zero in un numero finito di passi (al più \(n\), nel caso genericamente controllabile).

Questa è una differenza importante rispetto al tempo continuo:
nel continuo ``mandare un autovalore a \(-\infty\)'' non ha significato fisico,
mentre nel discreto
\[
z=0
\]
è un valore perfettamente lecito e corrisponde a un modo che si annulla in tempo finito.

\subsection{Interpretazione nel caso SISO}

Nel caso SISO, un sistema con tutti i poli in zero ha funzione di trasferimento del tipo
\[
G(z)=\frac{N(z)}{z^n}.
\]

Questa è la struttura tipica di un \textbf{filtro FIR} (\emph{Finite Impulse Response}),
cioè un sistema la cui risposta all'impulso si esaurisce in un tempo finito.

Per questo motivo un controllo dead-beat è molto rapido, ma presenta un problema pratico importante:
può richiedere azioni di controllo molto intense, quindi
\begin{itemize}
    \item si rischia di superare i limiti dell'attuatore;
    \item si può andare in saturazione;
    \item si perde la validità del modello lineare.
\end{itemize}

Quindi il dead-beat controller è concettualmente molto elegante,
ma spesso è troppo aggressivo per l'implementazione fisica.

\section{Sistemi stabilizzabili}

Si passa poi al caso in cui la coppia \((A,B)\) non sia completamente controllabile.

Consideriamo il sistema continuo
\[
\dot x=Ax+Bu.
\]

Se il sistema non è completamente controllabile, esiste un cambio di coordinate che separa
la parte controllabile da quella non controllabile:
\[
\dot z=
\begin{bmatrix}
A_c & A_{12}\\
0 & A_{nc}
\end{bmatrix}z
+
\begin{bmatrix}
B_c\\
0
\end{bmatrix}u,
\]
dove:
\begin{itemize}
    \item \(A_c\) descrive la parte controllabile;
    \item \(A_{nc}\) descrive la parte non controllabile.
\end{itemize}

Applicando una retroazione di stato
\[
u=-Kz,
\]
si può agire solo sulla parte controllabile, cioè sugli autovalori di
\[
A_c-B_cK_c.
\]

Gli autovalori della parte non controllabile, cioè quelli di \(A_{nc}\), \textbf{non possono essere spostati}.

\subsection{Definizione di stabilizzabilità}

Un sistema non completamente controllabile si dice \textbf{stabilizzabile} se tutti i modi non controllabili sono già stabili.

Nel caso tempo continuo:
\[
(A,B)\ \text{stabilizzabile}
\quad\Longleftrightarrow\quad
\Re(\lambda)<0\ \text{per tutti gli autovalori non controllabili}.
\]

Nel caso tempo discreto:
\[
(A,B)\ \text{stabilizzabile}
\quad\Longleftrightarrow\quad
|\lambda|<1\ \text{per tutti gli autovalori non raggiungibili/non controllabili}.
\]

Quindi:
\begin{itemize}
    \item la controllabilità completa permette di assegnare \emph{tutti} gli autovalori;
    \item la stabilizzabilità richiede solo che quelli non assegnabili siano già stabili.
\end{itemize}

\section{Perché la sola retroazione di stato non basta per il tracking perfetto}

Riprendiamo ora un sistema continuo
\[
\begin{cases}
\dot x=Ax+Bu,\\
y=Cx,
\end{cases}
\]
e consideriamo una legge di controllo del tipo
\[
u=-Kx+\bar r,
\]
dove \(\bar r\) rappresenta un riferimento costante.

Allora la dinamica diventa
\[
\dot x=(A-BK)x+B\bar r.
\]

All'equilibrio, imponendo \(\dot x=0\), si ottiene
\[
(A-BK)\bar x+B\bar r=0.
\]

Se la matrice
\[
N:=A-BK
\]
è invertibile, allora
\[
\bar x=-N^{-1}B\bar r.
\]

L'uscita a regime vale quindi
\[
\bar y=C\bar x=-CN^{-1}B\bar r.
\]

In generale \(\bar y\neq \bar r\), quindi la sola retroazione di stato \(-Kx\) non garantisce,
da sola, errore nullo a regime per un riferimento costante.

\section{Introduzione dell'integratore dell'errore}

Per eliminare l'errore a regime si introduce uno \textbf{stato integratore}.

Definiamo l'errore
\[
e(t)=r(t)-y(t)=r(t)-Cx(t),
\]
e introduciamo una nuova variabile di stato \(x_i(t)\) tale che
\[
\dot x_i(t)=e(t)=r(t)-Cx(t).
\]

Negli appunti compare prima la struttura in cui l'uscita dell'integratore entra come nuovo segnale di comando:
\[
u(t)=-Kx(t)+x_i(t).
\]

Sostituendo nella dinamica del sistema,
\[
\dot x=(A-BK)x+Bx_i.
\]

Quindi il sistema aumentato, con stato
\[
w=
\begin{bmatrix}
x\\
x_i
\end{bmatrix},
\]
si scrive
\[
\dot w
=
\begin{bmatrix}
A-BK & B\\
-C & 0
\end{bmatrix}
w
+
\begin{bmatrix}
0\\
I
\end{bmatrix}r.
\]

\subsection{Versione con ulteriore feedback sullo stato integratore}

Negli appunti compare poi una generalizzazione in cui si aggiunge un termine di retroazione anche sullo stato integratore:
\[
\dot x_i=r-Cx-K_i x_i.
\]

In questo caso il sistema aumentato diventa
\[
\dot w
=
\begin{bmatrix}
A-BK & B\\
-C & -K_i
\end{bmatrix}
w
+
\begin{bmatrix}
0\\
I
\end{bmatrix}r.
\]

L'idea è che anche la dinamica dell'integratore può essere modellata e ``chiusa''
attraverso un opportuno guadagno \(K_i\), così da permettere l'assegnamento
degli autovalori della matrice aumentata.

\section{Interpretazione dello stato aumentato}

L'introduzione di \(x_i\) ha due scopi:
\begin{itemize}
    \item memorizzare l'errore accumulato;
    \item forzare il sistema a correggere eventuali offset a regime.
\end{itemize}

Dal punto di vista del progetto, questo significa passare da un sistema di ordine \(n\)
a un sistema di ordine \(n+\ell\) (oppure \(n+1\) nel caso SISO),
su cui si può poi tentare una nuova retroazione di stato.

\section{Problema di osservazione dello stato}

La seconda grande parte della lezione introduce il problema dell'osservazione.

Finora si è sempre ipotizzato di conoscere tutto lo stato \(x(t)\),
ma nella pratica spesso sono misurabili solo:
\begin{itemize}
    \item l'ingresso \(u(t)\);
    \item l'uscita \(y(t)\).
\end{itemize}

Il sistema è
\[
\begin{cases}
\dot x=Ax+Bu,\\
y=Cx.
\end{cases}
\]

La domanda è:

\begin{quote}
\emph{Se lo stato non è direttamente accessibile, cosa posso ricostruire di \(x(t)\)
a partire da \(u(t)\), \(y(t)\) e dal modello \((A,B,C)\)?}
\end{quote}

\section{Primo tentativo: copia del modello}

Un primo tentativo consiste nel costruire una copia del sistema:
\[
\dot{\hat x}=A\hat x+Bu,
\]
dove \(\hat x(t)\) è una stima dello stato vero.

Se conoscessi esattamente:
\begin{itemize}
    \item il modello \(A,B\);
    \item la condizione iniziale \(x(0)\);
\end{itemize}
e scegliessi
\[
\hat x(0)=x(0),
\]
allora la stima coinciderebbe sempre con lo stato reale.

Infatti, definendo l'errore di stima
\[
e(t)=x(t)-\hat x(t),
\]
si ottiene
\[
\dot e=Ae.
\]

Se inizialmente \(e(0)=0\), allora \(e(t)\equiv 0\).

\subsection*{Problema pratico}

Questa idea però non è sufficiente nella realtà, perché:
\begin{itemize}
    \item non conosco esattamente la condizione iniziale;
    \item il modello può non essere perfetto;
    \item rumore e incertezze possono rendere l'errore non nullo.
\end{itemize}

Se ad esempio si usano matrici stimate \(\hat A,\hat B\), l'errore non soddisfa più una dinamica autonoma semplice,
e in generale non si può garantire che tenda a zero.

\section{Osservatore di Luenberger: idea}

Per correggere la stima si introduce una retroazione aggiuntiva basata sull'errore di uscita.

Poiché
\[
y=Cx,
\qquad
\hat y=C\hat x,
\]
la differenza
\[
y-\hat y = Cx-C\hat x = C(x-\hat x)=Ce
\]
contiene informazione sull'errore di stima.

Si costruisce allora un osservatore del tipo
\[
\dot{\hat x}=A\hat x+Bu+L\bigl(y-\hat y\bigr),
\]
cioè
\[
\dot{\hat x}=A\hat x+Bu+L\bigl(y-C\hat x\bigr),
\]
dove \(L\) è il guadagno dell'osservatore.

Negli appunti la stessa idea compare anche nella forma equivalente
\[
\dot{\hat x}=A\hat x+Bu-\mu(C\hat x-y),
\]
dove il simbolo del guadagno può cambiare, ma il significato resta lo stesso.

\section{Dinamica dell'errore di osservazione}

Definendo ancora
\[
e=x-\hat x,
\]
si ha
\[
\dot e
=
(Ax+Bu)-\bigl(A\hat x+Bu+L(y-C\hat x)\bigr).
\]

Poiché \(y=Cx\),
\[
\dot e=A(x-\hat x)-LC(x-\hat x),
\]
quindi
\[
\dot e=(A-LC)e.
\]

Questo è il punto chiave:
scegliendo \(L\), posso modificare gli autovalori della matrice
\[
A-LC.
\]

Quindi il problema dell'osservazione è duale a quello della retroazione di stato:
\begin{itemize}
    \item nella retroazione di stato si progettava \(K\) per agire su \(A-BK\);
    \item nell'osservatore si progetta \(L\) per agire su \(A-LC\).
\end{itemize}

\section{Commento finale sulla lezione}

La lezione mette insieme tre idee importanti:

\begin{enumerate}
    \item nel discreto è possibile progettare un \emph{dead-beat controller},
    assegnando tutti gli autovalori in zero;

    \item anche se un sistema non è completamente controllabile,
    può comunque essere \emph{stabilizzabile}, purché i modi non controllabili siano già stabili;

    \item quando la sola retroazione di stato non basta a garantire tracking perfetto,
    si introduce uno stato integratore e si passa a un sistema aumentato;

    \item quando lo stato non è interamente misurabile,
    nasce il problema dell'osservazione, che porta all'introduzione dell'osservatore dinamico.
\end{enumerate}

\section{Sintesi della lezione}

I punti principali sono:
\begin{itemize}
    \item per un sistema discreto
    \[
    x(k+1)=Ax(k)+Bu(k),
    \qquad
    u(k)=-Kx(k),
    \]
    è possibile assegnare gli autovalori del sistema chiuso in
    \[
    z=0,
    \]
    ottenendo un \textbf{dead-beat controller};

    \item il dead-beat controllo è molto rapido, ma può richiedere azioni di controllo troppo elevate;

    \item un sistema non completamente controllabile è \textbf{stabilizzabile}
    se i modi non controllabili sono già stabili;

    \item per eliminare l'errore a regime si introduce un integratore dell'errore,
    ottenendo un sistema aumentato;

    \item se lo stato non è direttamente misurabile, si introduce una stima \(\hat x\);

    \item una semplice copia del modello non basta in presenza di incertezze;

    \item usando l'errore di uscita \(y-\hat y\) si costruisce un osservatore dinamico:
    \[
    \dot{\hat x}=A\hat x+Bu+L(y-C\hat x);
    \]

    \item l'errore di osservazione evolve secondo
    \[
    \dot e=(A-LC)e,
    \]
    e quindi il progetto dell'osservatore consiste nello scegliere \(L\)
    in modo che questa dinamica sia stabile e sufficientemente veloce.
\end{itemize}